\chapter*{Abstract}

The rapid growth of digital images and video, together with stricter and evolving regulations on personal data protection, has increased the need for reliable methods to safeguard sensitive information. Manual redaction is time-consuming and difficult to scale, while fully automated approaches often lack consistent accuracy in complex visual environments. This thesis presents \gsr{SafeVision}, a system for detecting and anonymizing sensitive content in visual media. 

The system combines \gls{object detection} and text analysis to identify sensitive content in visual media. It detects faces, text, and general objects such as people and vehicles. In addition, custom models were trained and integrated into the system to classify and label domain specific entities. 

The system is designed as a scalable and user centered architecture that supports reliable redaction. The machine learning components are integrated into a Python backend built with FastAPI for asynchronous processing, alongside a web based frontend developed with React and TypeScript. Users can review and adjust detections before redaction, enabling user guided refinement that improves accuracy and usability. The modular design, with clear separation between components, supports extensibility and maintainability, while client side storage reduces data transfer and enhances privacy.